% TRACE-SOURCE-ID: TEX-LAYOUT-20260819-V1
% TRACE-PURPOSE: two-column order, caption, footnote and citation fidelity
\documentclass[10pt,a4paper]{ctexart}
\usepackage[margin=18mm]{geometry}
\usepackage{amsmath,fancyhdr,hyperref,multicol,xcolor}
\hypersetup{
  pdftitle={Traceable Layout Evaluation},
  pdfauthor={Document Parser QA},
  pdfsubject={TEX-LAYOUT-20260819-V1},
  pdfkeywords={TRACE-L-DOC-001, two-column, caption, footnote}
}
\pagestyle{fancy}
\fancyhf{}
\fancyhead[L]{TRACE-RUNNING-HEADER-L}
\fancyhead[R]{TEX-LAYOUT-20260819-V1}
\fancyfoot[C]{Page token PL-42 | \thepage}

\begin{document}
\begin{center}
{\LARGE 双栏与对象关系测试}\par
{\small TRACE-L-DOC-001 | Source: TEX-LAYOUT-20260819-V1}
\end{center}

\begin{multicols}{2}
\section{左栏起点 TRACE-L-COL-A01}
阅读顺序声明为 READ-A-B-C-D。A 段必须先于 B 段，随后进入右栏的 C 段与 D 段。
本段包含条件 $\alpha \leq \beta + 0.05$，不等号与小数必须保留。

\subsection{对象说明}
下面的框模拟独立图形对象。它不携带外部资源，但图注与对象必须保持邻接关系。

\begin{center}
\fbox{\parbox{0.82\columnwidth}{
\centering
TRACE-L-FIG-101\\[3pt]
Figure checksum 7F3A-91C2\\
Flow: INPUT $\rightarrow$ VERIFY $\rightarrow$ OUTPUT
}}

\small 图 1：验证流程 TRACE-L-CAP-111
\end{center}

对象校验使用固定 checksum，不得把字符 F、A、C 与数字混淆。

\columnbreak
\section{右栏起点 TRACE-L-COL-B01}
C 段延续阅读顺序。这里包含脚注引用\footnote{TRACE-L-FOOT-121：脚注事实为 footnote-value-19.25。}，
脚注标记应与脚注正文建立关系，而不是并入相邻标题。

\subsection{约束公式}
\begin{equation}
\alpha \leq \beta + 0.05
\tag{TRACE-L-EQ-131}
\end{equation}

D 段完成 READ-A-B-C-D。页面底部的 Page token PL-42 是重复页脚事实，
可标记为版面噪声，但不得擅自改写其字符。
\end{multicols}

\newpage
\section{引用与附录 TRACE-L-REF-201}
正式引用：Doe, J. (2026), Traceable Layout Recovery, DOI 10.5555/trace.2026.0819。

\begin{quote}
图、图注、脚注、公式与引用必须通过已有对象和来源证据建立关系。
\end{quote}

本页再次出现运行页眉与 Page token PL-42，用于检测重复页眉页脚的解析行为。

\end{document}
